% Component diagram: three executables over one Qt-free core; strict downward layering.
% Consumed by the system-overview section. Bare tikzpicture; the section wraps it in a figure.
\begin{tikzpicture}[node distance=7mm and 8mm]

  % --- executables (top) ---
  \node[consumer,text width=40mm] (gui)
    {\textbf{qtrocket} (GUI)\\[1pt]\footnotesize\ttfamily gui/ GuiRunner, MainWindow\\
     \rmfamily Qt6 Widgets \textbullet\ qcustomplot\\\footnotesize the only Qt-dependent layer};
  \node[consumer,text width=38mm,right=of gui] (cli)
    {\textbf{qtrocket-cli} (REPL)\\[1pt]\footnotesize\ttfamily cli/ Repl\\
     \rmfamily headless, Qt-free\\\footnotesize scriptable via stdin};
  \node[consumer,text width=38mm,right=of cli] (viz)
    {\textbf{visualizer} (OpenGL)\\[1pt]\footnotesize\ttfamily visualizer/ RocketMesh\\
     \rmfamily standalone 3D viewer\\\footnotesize reads the same part tree};

  % --- controller / core lib ---
  \node[resolver,text width=104mm,below=11mm of cli] (core)
    {qtrocket\_core \textcolor{qrGray}{(static lib)}\\[2pt]
     \normalfont\footnotesize QtRocket controller singleton: owns (RocketModel, Propagator) pair,\\
     Environment, MotorModelDatabase \textbullet\ launchRocket() / setInitialState() / getStates()};

  % --- model + sim (side by side) ---
  \node[comp,text width=50mm,below=11mm of core.south west,anchor=north west] (model)
    {\textbf{model/} \footnotesize physical description\\[1pt]
     \footnotesize Part composite tree \textbullet\ placement resolver\\
     MotorModel + ThrustCurve \textbullet\ motor DB\\
     RSE / RASP / thrustcurve.org ingest\\
     DesignSerializer (.qrd)};
  \node[comp,text width=50mm,below=11mm of core.south east,anchor=north east] (sim)
    {\textbf{sim/} \footnotesize numerics\\[1pt]
     \footnotesize Propagator ODE loop \textbullet\ RK4 / RKF45\\
     Environment: gravity, atmosphere,\\
     geoid, wind \textbullet\ Aero value types\\
     TrajectoryStatistics};

  % --- utils (foundation) ---
  % anchored off the model box (the taller of the two mid-layer boxes), not a fixed
  % distance from core, so the box can never paint over the layer above it.
  \node[comp,text width=104mm,below=7mm of model.south west,anchor=north west,fill=qrGray!8,draw=qrGray] (utils)
    {\textbf{utils/} \footnotesize Logger singleton \textbullet\ math typedefs over Eigen
     (\texttt{Vector3}, \texttt{Matrix3}, \texttt{Quaternion}) \textbullet\ CurlConnection (libcurl)\\
     \footnotesize never includes model/ or sim/ --- keeps the dependency graph acyclic};

  % --- model <-> sim bridge note ---
  \node[draw=qrTeal,fill=qrTeal!7,rounded corners=2pt,line width=0.5pt,
        text width=36mm,font=\footnotesize,right=8mm of core.east,anchor=west] (bridge)
    {\textbf{Propagatable} is the model$\leftrightarrow$sim bridge: \texttt{getForces()},
     \texttt{getTorques()}, \texttt{getMass(t)}, \texttt{terminateCondition()}.};

  % --- edges ---
  \draw[uses] (gui.south) -- (core.north -| gui.south) node[flowlbl,midway]{links};
  \draw[uses] (cli.south) -- (core.north);
  \draw[uses] (viz.south) -- (core.north -| viz.south) node[flowlbl,midway]{links};
  \draw[dep] (core.south -| model.north) -- (model.north);
  \draw[dep] (core.south -| sim.north) -- (sim.north);
  \draw[dep] (model.south) -- (model.south |- utils.north);
  \draw[dep] (sim.south) -- (sim.south |- utils.north);
  \draw[->,qrTeal,line width=0.6pt,dashed] (bridge.west) -- (core.east);

\end{tikzpicture}
